\section{Introducción}
\label{sec:intro}
En los últimos años, hemos visto un aumento considerable en el uso de la IA para suplantar la identidad de alguien de forma maliciosa clonando su voz. Estos son los conocidos \emph{deepfakes}. El reconocimiento de dichos deepfakes es de vital importancia, pues suponen una gran amenaza. Por otro lado, las técnicas de detección de \emph{deepfakes} no han avanzado al mismo ritmo, y esto lo vemos en la gran cantidad de estafas que se producen de este tipo.

\hspace{.7cm} 
Además, las personas ya no somos capaces de diferenciar entre una voz real y una generada, como se ve en \cite{1Barrington2025, 2Nadine2025}, donde hacen sendos estudios exponiendo a personas reales a diversos audios, algunos reales, algunos generados y algunos que clonan voces reales. En dichos estudios, se ve que los que son 100\% generados nos cuesta poco detectarlos, pero entre los reales y los clones no diferenciamos bien, y rondamos el 60\% de precisión clasificando dichas voces en reales o falsas. Por esto, no podemos fiarnos más de nuestras capacidades a la hora de diferenciar cuando una voz es real o no.

\hspace{.7cm}	
En este texto, se proponen varias técnicas que podrían ser útiles para la detección de audios generados por IA. Para esta tarea es fundamental sacar las características de la voz, así como gráficas, para poder entrenar los modelos con dichos datos. Entre los audios debe haber tanto reales como generados, y estar medianamente equilibrados, para que los modelos aprendan las diferencias, y así sepan clasificar un audio que nunca han visto.

\hspace{.7cm}
Pero el rendimiento del modelo de clasificación no es lo único que nos interesa. Hoy en día, los modelos de IA se han convertido en \emph{cajas negras}, donde no sabemos cómo hace las predicciones, dando lugar a problemas de confianza de los humanos respecto de dichas predicciones, como es en el ámbito médico o ingenieril, donde las decisiones que tomen las IAs deben ser precisas, puesto que la vida de las personas puede estar en juego. Para solventar este problema, ha surgido la IA explicable (XAI), la cual toma dichos modelos \emph{caja negra} y da una explicación de qué características de los datos afectan más a la decisión del modelo. En esta investigación, vamos a aplicar XAI para ver qué afecta más a la clasificación de los audios dentro de los distintos grupos que mencionamos antes.